\begin{recipe}
[
    preparationtime = {\unit[20--30]{min}},
    portion = {\portion{3}},
    calory = {260 kcal | 5g Protein | 54g KH | 3g Fett},
    rating = {2},
    source = {Adrian Hieber},
    price = {0,73€}
]
{Schoko-Erdbeeren Mochi}

\introduction{%
    Weich, schokoladig und leicht fruchtig durch die Erdbeeren. 
    Perfekt als kleines Dessert oder Snack.
}

\ingredients{%
    100 g & Klebreismehl\\
    25 g & Kakaopulver\\
    35 g & Brauner Zucker\\
    175 ml & Wasser\\
    9 & Erdbeeren (klein-mittel)\\
    Etwas & Kartoffelstärke zum Bestäuben
}

\preparation{%
    \begin{enumerate}
        \item Mehl, Zucker und Kakao vermischen.
        \item Mit Wasser zu einem glatten Teig rühren.
        \item Zubereiten:
        \begin{enumerate}
            \renewcommand{\theenumi}{\arabic{enumi}}
            \renewcommand{\theenumii}{\theenumi.\arabic{enumii}}
            \item \textbf{Pfanne:} In einer beschichteten Pfanne bei mittlerer Hitze unter ständigem Rühren erhitzen, bis eine zähe Masse entsteht (ca. 6 Minuten).
            \item \textbf{Mikrowelle:} In der Mikrowelle 5 Minuten bei hoher Leistung garen und zwischendurch umrühren.
        \end{enumerate}
        \item In 9 Teile schneiden, Arbeitsfläche mit Stärke bestäuben.
        \item Erdbeeren trocken tupfen, jeweils ein Stück Teig darum legen und zu Kugeln formen.
    \end{enumerate}
}

\suggestion[Variation]{%
    Funktioniert auch sehr gut mit Himbeeren oder einem kleinen Stück dunkler Schokolade als Kern.
}

\hint{%
    Erdbeeren definitiv gut abtrocknen, sonst löst sich der Teig auf, gefroene Früchte sind deshalb auch problematisch.
    Klebreismehl kaufen und nicht selber herstellen, das funktioniert nicht, ich spreche aus Erfahrung.
}

\end{recipe}